\chapter{Laboratorio Python: Data Understanding e Data Preparation}
\label{chap:dm-python-lab-du-dp}

In questo capitolo viene presentata la trasposizione computazionale delle metodologie teoriche di Data Understanding (DU) e Data Preparation (DP), facendo riferimento all'attività di laboratorio svolta con lo stack scientifico di Python (\texttt{numpy}, \texttt{pandas}, \texttt{scipy}, \texttt{matplotlib}, \texttt{seaborn} e \texttt{scikit-learn}).

Il laboratorio impiega il dataset \textbf{Tips} (arricchito con la tabella \textbf{Tips\_Age}), relativo alle abitudini di spesa e di mancia dei clienti di un ristorante, per illustrare l'intero flusso operativo: dall'ingestione e integrazione dei dati, alla diagnosi e imputazione dei valori mancanti, fino all'ingegnerizzazione delle feature, normalizzazione e riduzione dimensionale con PCA e $t$-SNE.

\begin{noteblock}{Ambiente Interattivo e Riproducibilità del Laboratorio}
    Tutti gli esempi, le pipeline di trasformazione e i grafici presentati in questo capitolo sono riproducibili ed eseguibili interattivamente tramite i notebook Jupyter inclusi nella cartella \texttt{Data-Mining/notebooks/}.
    
    \textbf{Avvio rapido con Docker}:
    \begin{itemize}[noitemsep]
        \item Da terminale (nella radice del progetto): \texttt{docker compose up jupyter}
        \item Oppure con doppio clic su Windows: \texttt{scripts/start-jupyter.cmd}
        \item Interfaccia web: \href{http://localhost:8888/lab}{\texttt{http://localhost:8888/lab}} (senza password né token).
    \end{itemize}
    Per eseguire direttamente gli script senza interfaccia grafica da riga di comando:
    \begin{itemize}[noitemsep]
        \item \texttt{docker compose run --rm app python3 Data-Mining/notebooks/generate\_dispensa\_figures.py}
    \end{itemize}
\end{noteblock}

% ====================================================================
% SEZIONE 5.1: AMBIENTE, INGESTIONE E DATA INTEGRATION
% ====================================================================
\section{Ambiente di Lavoro, Ingestione e Data Integration}
\label{sec:dm-lab-ingestion-integration}

\begin{remark}[Riferimento al Notebook Interattivo]
    Il codice delle sezioni \ref{sec:dm-lab-ingestion-integration}, \ref{sec:dm-lab-missing-values} e \ref{sec:dm-lab-multivariate-exploration} corrisponde al notebook Jupyter interattivo \href{http://localhost:8888/lab/tree/notebooks/1_basics_and_understanding.ipynb}{\textbf{Notebook 1 (Data Understanding)}} (file \texttt{1\_basics\_and\_} \texttt{understanding.ipynb}).
\end{remark}

Lo stack software fondamentale per l'analisi esplorativa dei dati in Python comprende le seguenti librerie:

\begin{lstlisting}[language=Python, style=code]
import numpy as np
import pandas as pd
import scipy.stats as stats
import matplotlib.pyplot as plt
import seaborn as sns
from sklearn.impute import SimpleImputer, KNNImputer
from sklearn.preprocessing import StandardScaler, MinMaxScaler, KBinsDiscretizer, OneHotEncoder, TargetEncoder
from sklearn.decomposition import PCA
from sklearn.manifold import TSNE
\end{lstlisting}

\subsection{Caricamento e Ispezione Strutturale}

Il dataset principale viene acquisito tramite file delimitato da tabulazione (\texttt{tsv}) impostando l'indice sulla prima colonna:

\begin{lstlisting}[language=Python, style=code]
tips = pd.read_csv('tips.csv', sep='\t', index_col=0)
print(tips.head())
tips.info()
\end{lstlisting}

Il dataset contiene inizialmente $244$ osservazioni e $7$ colonne: \texttt{total\_bill} (importo totale del conto), \texttt{tip} (mancia), \texttt{sex} (genere del pagante), \texttt{smoker} (presenza di fumatori al tavolo), \texttt{day} (giorno della settimana), \texttt{time} (pranzo o cena), \texttt{size} (numero di coperti al tavolo).

\subsection{Data Integration e Join Relazionale}

Nel laboratorio, l'età dei clienti è memorizzata in una sorgente secondaria (\texttt{tips\_age.csv}). L'integrazione delle due tabelle avviene tramite una join sull'identificativo primario \texttt{ID}:

\begin{lstlisting}[language=Python, style=code]
tips_age = pd.read_csv('tips_age.csv', sep='\t', index_col=0)
df = tips.join(tips_age, on='ID')
# Rimozione di eventuali colonne ridondanti generate dalla sorgente
df = df.drop(columns=['gender'], errors='ignore')
\end{lstlisting}

\subsection{Rilevamento e Rimozione di Record Duplicati}

L'ispezione delle righe duplicate richiede la verifica dell'uguaglianza tra tutti gli attributi del record:

\begin{lstlisting}[language=Python, style=code]
# Visualizzazione di tutti i duplicati (tutte le copie)
duplicated_rows = df[df.duplicated(keep=False)]
print(f"Righe duplicate rilevate: {len(duplicated_rows)}")

# Rimozione mantenendo la prima occorrenza
df = df.drop_duplicates()
\end{lstlisting}

\subsection{Bonifica di Valori Non Validi e Cast di Tipo}

Spesso i dati contengono errori di trascrizione o stringhe non valide inserite al posto di campi numerici (ad esempio il valore fittizio \texttt{'aaaaa'} nella colonna \texttt{total\_bill}). Ciò costringe Pandas a tipizzare la colonna come generico tipo \texttt{object}. La bonifica avviene mediante sostituzione puntuale e conversione esplicita di tipo:

\begin{lstlisting}[language=Python, style=code]
# Sostituzione della stringa anomala con zero o NaN
df["total_bill"] = df["total_bill"].replace("aaaaa", np.nan)

# Cast esplicito a float numerico
df["total_bill"] = pd.to_numeric(df["total_bill"])
\end{lstlisting}

% ====================================================================
% SEZIONE 5.2: DIAGNOSI E IMPUTAZIONE DEI VALORI MANCANTI
% ====================================================================
\section{Diagnosi e Imputazione dei Valori Mancanti}
\label{sec:dm-lab-missing-values}

\subsection{Audit dei Missing Values}

Il conteggio dei valori mancanti (\texttt{NaN}) per ciascun attributo si ottiene combinando \texttt{isnull()} e \texttt{sum()}:

\begin{lstlisting}[language=Python, style=code]
# Riepilogo dei valori nulli per colonna
null_counts = df.isnull().sum()
print(null_counts[null_counts > 0])

# Selezione dei soli record incompleti
incomplete_records = df[df.isnull().any(axis=1)]
\end{lstlisting}

\subsection{Strategie di Rimozione vs Imputazione Statistica}

\begin{itemize}
    \item \textbf{Listwise Deletion}: Eliminazione integrale di tutti i record con almeno un valore nullo:
\begin{lstlisting}[language=Python, style=code]
df_cleaned = df.dropna()
\end{lstlisting}
    \item \textbf{Imputazione con la Mediana Condizionale}: Quando l'attributo \texttt{age} presenta valori mancanti o codificati con la convenzione sentinella $0$, l'imputazione mediante la mediana calcolata sull'intero dataset potrebbe appiattire la variabilità tra fasce orarie distinte (\textit{Lunch} vs \textit{Dinner}). Si esegue pertanto un'imputazione stratificata:
\begin{lstlisting}[language=Python, style=code]
df['age'] = df['age'].replace(0, np.nan)
# Imputa con la mediana di gruppo calcolata sulla fascia oraria (time)
df['age'] = df['age'].fillna(df.groupby('time')['age'].transform('median'))
\end{lstlisting}
\end{itemize}

\subsection{Imputazione Avanzata con Scikit-Learn}

Il modulo \texttt{sklearn.impute} mette a disposizione trasformatori integrabili nelle pipeline di machine learning:

\begin{lstlisting}[language=Python, style=code]
# 1. SimpleImputer per feature numeriche (media o mediana)
imp_mean = SimpleImputer(strategy='mean')
df['tip'] = imp_mean.fit_transform(df[['tip']])

# 2. SimpleImputer per feature categoriche (moda/frequenza massima)
imp_mode = SimpleImputer(strategy='most_frequent')
df['smoker'] = imp_mode.fit_transform(df[['smoker']]).ravel()

# 3. KNNImputer (imputazione multivariata basata su k-Nearest Neighbors)
knn_imp = KNNImputer(n_neighbors=2, weights='distance')
num_cols = ['total_bill', 'tip', 'size', 'age']
df[num_cols] = knn_imp.fit_transform(df[num_cols])
\end{lstlisting}

% ====================================================================
% SEZIONE 5.3: ESPLORAZIONE MULTIVARIATA E CORRELAZIONE
% ====================================================================
\section{Esplorazione Multivariata e Analisi di Correlazione}
\label{sec:dm-lab-multivariate-exploration}

\subsection{Tabelle di Contingenza e Distribuzioni Condizionate}

Per studiare le associazioni tra variabili categoriche si utilizzano tabelle di contingenza normalizzate (\textit{cross-tabulation}):

\begin{lstlisting}[language=Python, style=code]
# Tabella di frequenza congiunta tra genere e abitudine al fumo
smoker_xt = pd.crosstab(df['smoker'], df['sex'])

# Normalizzazione per riga (proporzioni percentuali)
smoker_pct = smoker_xt.div(smoker_xt.sum(1).astype(float), axis=0)
smoker_pct.plot(kind='bar', stacked=False, title='Genere per Categoria Fumatore')
plt.ylabel('Frequenza Relativa')
plt.show()
\end{lstlisting}

\subsection{Matrici di Correlazione e Heatmap}

Pandas consente di stimare simultaneamente i coefficienti di correlazione lineare (Pearson) e di rango (Spearman, Kendall) sui soli attributi numerici:

\begin{lstlisting}[language=Python, style=code]
# Calcolo della matrice di correlazione di Pearson
corr_pearson = df.corr(numeric_only=True, method='pearson')

# Calcolo con il coefficiente di Spearman (monotonia robusta)
corr_spearman = df.corr(numeric_only=True, method='spearman')

# Visualizzazione grafica tramite heatmap con annotazioni numeriche
plt.figure(figsize=(8, 6))
sns.heatmap(corr_pearson, annot=True, cmap='coolwarm', fmt='.2f', vmin=-1, vmax=1)
plt.title('Matrice di Correlazione di Pearson (Dataset Tips)')
plt.show()
\end{lstlisting}

\subsection{Scatter Matrix e Pairplot}

L'ispezione visiva incrociata delle relazioni bivariate si realizza tramite \texttt{scatter\_matrix} o Seaborn \texttt{pairplot}:

\begin{lstlisting}[language=Python, style=code]
# Scatter Matrix di Pandas
pd.plotting.scatter_matrix(df[['total_bill', 'tip', 'size', 'age']], figsize=(9, 9), diagonal='kde')
plt.show()

# Pairplot con codice colore basato sul momento del pasto (time)
sns.pairplot(df, vars=['total_bill', 'tip', 'size'], hue='time', palette='Dark2')
plt.show()
\end{lstlisting}

\subsection{Feature Engineering: Costruzione di Indicatori di Dominio}

La conoscenza del dominio consente di combinare variabili grezze per sintetizzare nuovi attributi ad alto valore interpretativo:

\begin{lstlisting}[language=Python, style=code]
# 1. Percentuale effettiva di mancia rispetto al conto
df['tip_ratio'] = df['tip'] / df['total_bill']

# 2. Spesa media per persona (Bill Per Person - BPP)
df['bpp'] = df['total_bill'] / df['size']

# Ispezione con Boxplot condizionato per genere
df.boxplot(column=['bpp'], by='sex')
plt.title('Spesa Media per Persona per Genere')
plt.suptitle('')
plt.show()
\end{lstlisting}

% ====================================================================
% SEZIONE 5.4: FEATURE TRANSFORMATION E ENCODING
% ====================================================================
\section{Feature Transformation, Discretizzazione e Encoding}
\label{sec:dm-lab-transformation-encoding}

\begin{remark}[Riferimento al Notebook Interattivo]
    Il codice delle sezioni \ref{sec:dm-lab-transformation-encoding} e \ref{sec:dm-lab-dimensionality-reduction} corrisponde al notebook Jupyter interattivo \href{http://localhost:8888/lab/tree/notebooks/2_feature_engineering_and_data_representation.ipynb}{\textbf{Notebook 2 (Feature Engineering)}}:
    \begin{center}
        \small \href{http://localhost:8888/lab/tree/notebooks/2_feature_engineering_and_data_representation.ipynb}{\texttt{2\_feature\_engineering\_and\_data\_representation.ipynb}}
    \end{center}
\end{remark}

\subsection{Discretizzazione con KBinsDiscretizer}

La libreria Scikit-Learn fornisce \texttt{KBinsDiscretizer} per convertire variabili continue in classi ordinali discrete secondo tre strategie:

\begin{lstlisting}[language=Python, style=code]
# Discretizzazione a 4 intervalli con strategie diverse
disc_uniform = KBinsDiscretizer(n_bins=4, encode='ordinal', strategy='uniform')    # Equal-Width
disc_quantile = KBinsDiscretizer(n_bins=4, encode='ordinal', strategy='quantile')  # Equal-Frequency
disc_kmeans = KBinsDiscretizer(n_bins=4, encode='ordinal', strategy='kmeans')      # Basata su Cluster

df['total_bill_bins'] = disc_uniform.fit_transform(df[['total_bill']])
print(df['total_bill_bins'].value_counts())
\end{lstlisting}

\subsection{Normalizzazione dei Dati: StandardScaler vs MinMaxScaler}

\begin{itemize}
    \item \textbf{StandardScaler}: Trasforma la feature per avere media zero e deviazione standard unitaria ($\mu = 0, \sigma = 1$).
    \item \textbf{MinMaxScaler}: Riscala linearmente l'intervallo originario entro $[0, 1]$.
\end{itemize}

\begin{lstlisting}[language=Python, style=code]
num_cols = ['total_bill', 'tip', 'age']

# Standardizzazione Z-Score
scaler_std = StandardScaler()
df_standardized = df.copy()
df_standardized[num_cols] = scaler_std.fit_transform(df[num_cols])

# Normalizzazione Min-Max
scaler_minmax = MinMaxScaler()
df_normalized = df.copy()
df_normalized[num_cols] = scaler_minmax.fit_transform(df[num_cols])

# Ripristino dei valori originari tramite trasformazione inversa
df_recovered = scaler_minmax.inverse_transform(df_normalized[num_cols])
\end{lstlisting}

\subsection{Codifica di Variabili Categoriche: One-Hot vs Target Encoding}

\begin{lstlisting}[language=Python, style=code]
# 1. One-Hot Encoding (evita di creare ordinalita fittizia)
ohe = OneHotEncoder(handle_unknown='ignore', sparse_output=False)
encoded_days = ohe.fit_transform(df[['day']])
encoded_df = pd.DataFrame(encoded_days, columns=ohe.get_feature_names_out(['day']))

# 2. Target Encoding (sostituisce la categoria con la media della target)
target_enc = TargetEncoder(smooth='auto')
# Codifica 'day' in base all'importo medio della mancia 'tip'
df['day_target_encoded'] = target_enc.fit_transform(df[['day']], df['tip'])
\end{lstlisting}

% ====================================================================
% SEZIONE 5.5: RIDUZIONE DIMENSIONALE: PCA E T-SNE
% ====================================================================
\section{Riduzione Dimensionale: PCA e $t$-SNE}
\label{sec:dm-lab-dimensionality-reduction}

\subsection{Principal Component Analysis (PCA)}

Prima di applicare la PCA è indispensabile standardizzare le feature continue affinché attributi con varianze ampie non dominino la decomposizione spettrale:

\begin{lstlisting}[language=Python, style=code]
# Selezione attributi quantitativi
X = df[['total_bill', 'tip', 'size', 'age', 'tip_ratio', 'bpp']].dropna()
X_scaled = StandardScaler().fit_transform(X)

# Esecuzione della PCA
pca = PCA()
X_pca = pca.fit_transform(X_scaled)

# Autovalori (quota di varianza spiegata per componente)
eigenvalues = pca.explained_variance_
# Percentuale di varianza spiegata
explained_var_ratio = pca.explained_variance_ratio_
# Varianza cumulata
cumulative_var = np.cumsum(explained_var_ratio)
\end{lstlisting}

\subsubsection{Costruzione dello Scree Plot}
Lo Scree Plot consente di individuare visivamente il "gomito" (\textit{elbow}) oltre il quale le componenti spiegano quote trascurabili di varianza:

\begin{lstlisting}[language=Python, style=code]
plt.figure(figsize=(8, 4))
plt.plot(range(1, len(eigenvalues) + 1), eigenvalues, 'o-', color='navy', label='Autovalori singoli')
plt.title('Scree Plot degli Autovalori (PCA)')
plt.xlabel('Numero Componente Principale')
plt.ylabel('Autovalore (Varianza spiegata)')
plt.grid(True, linestyle='--', alpha=0.5)
plt.legend()
plt.show()
\end{lstlisting}

\subsection{Proiezione Non Lineare con $t$-SNE}

Il \textbf{$t$-Distributed Stochastic Neighbor Embedding ($t$-SNE)} è una tecnica non supervisionata non lineare specificamente progettata per preservare le strutture di vicinato locale nei dati ad alta dimensionalità:

\begin{lstlisting}[language=Python, style=code]
tsne = TSNE(
    n_components=2,
    perplexity=30,      # Equilibrio tra granularita locale e globale (tipicamente 5-50)
    learning_rate=200,  # Passo di ottimizzazione della discesa del gradiente
    max_iter=1000,      # Iterazioni di minimizzazione della divergenza KL
    random_state=42
)
X_tsne = tsne.fit_transform(X_scaled)
\end{lstlisting}

\subsection{Confronto Metodologico tra PCA e $t$-SNE}

\begin{table}[htbp]
    \centering
    \small
    \begin{tabularx}{\textwidth}{lXX}
        \toprule
        \textbf{Caratteristica} & \textbf{PCA (Principal Component Analysis)} & \textbf{$t$-SNE ($t$-Distributed SNE)} \\
        \midrule
        \textbf{Tipologia} & Lineare, globale, deterministica. & Non lineare, locale, probabilistica/stocastica. \\
        \addlinespace
        \textbf{Funzione Obiettivo} & Massimizza la varianza globale proiettata (decomposizione agli autovalori). & Minimizza la divergenza di Kullback-Leibler tra distribuzioni di probabilità di vicinato. \\
        \addlinespace
        \textbf{Interpretabilità} & Gli assi rappresentano combinazioni lineari degli attributi (pesi/loading noti). & Gli assi non hanno significato fisico; contano solo le distanze relative tra punti. \\
        \addlinespace
        \textbf{Distanze Lontane} & Preserva fedelmente le distanze euclidee su larga scala. & Non preserva le distanze globali tra cluster distanti (la grandezza dei cluster non riflette la densità originaria). \\
        \addlinespace
        \textbf{Uso Tipico} & Riduzione preliminare di feature per modelli predittivi, compressione, eliminazione di collinearità. & Visualizzazione esplorativa avanzata e separazione visiva di cluster complessi (2D/3D). \\
        \bottomrule
    \end{tabularx}
    \caption{Confronto metodologico tra PCA e $t$-SNE nell'analisi dei dati.}
    \label{tab:dm-pca-vs-tsne}
\end{table}

\begin{lstlisting}[language=Python, style=code]
# Visualizzazione comparata delle due proiezioni bidimensionali
fig, (ax1, ax2) = plt.subplots(1, 2, figsize=(14, 5))

# Scatter PCA
ax1.scatter(X_pca[:, 0], X_pca[:, 1], c='navy', alpha=0.7, edgecolors='k')
ax1.set_title('Proiezione Lineare PCA (PC1 vs PC2)')
ax1.set_xlabel('Prima Componente Principale')
ax1.set_ylabel('Seconda Componente Principale')
ax1.grid(True, linestyle='--', alpha=0.5)

# Scatter t-SNE
ax2.scatter(X_tsne[:, 0], X_tsne[:, 1], c='crimson', alpha=0.7, edgecolors='k')
ax2.set_title('Proiezione Non Lineare t-SNE')
ax2.set_xlabel('Dimensione t-SNE 1')
ax2.set_ylabel('Dimensione t-SNE 2')
ax2.grid(True, linestyle='--', alpha=0.5)

plt.tight_layout()
plt.show()
\end{lstlisting}
